% Generated draft: requires mathematical and source review.
\documentclass[11pt]{article}
\usepackage[margin=25mm]{geometry}
\usepackage{fontspec}
\setmainfont{Latin Modern Roman}
\usepackage{amsmath,amssymb,mathtools,unicode-math}
\setmathfont{Latin Modern Math}
\usepackage{enumitem,needspace,xurl,hyperref}
\setlength{\parindent}{0pt}
\setlength{\parskip}{5pt}
\setlength{\emergencystretch}{3em}
\newcommand{\N}{\mathbb{N}}
\newcommand{\Z}{\mathbb{Z}}
\newcommand{\Q}{\mathbb{Q}}
\newcommand{\R}{\mathbb{R}}
\newcommand{\C}{\mathbb{C}}
\newcommand{\F}{\mathbb{F}}
\newcommand{\A}{\mathbb{A}}
\newcommand{\PP}{\mathbb{P}}
\newcommand{\E}{\mathbb{E}}
\newcommand{\eps}{\varepsilon}
\newcommand{\dd}{\,\mathrm{d}}
\begin{document}
\section*{UnsolvedMath ID: 800261}
\textbf{Source problem code:} catalogue entry
\subsection{Definitions and mathematical statement}
For a positive integer \(n\), let
\[
SL(n,\C)=\{A\in\operatorname{Mat}_{n\times n}(\C):\det A=1\},
\]
viewed as a complex manifold of dimension \(n^2-1\). A holomorphic immersion into \(\C^{n^2-1}\) is a holomorphic map
\[
f:SL(n,\C)\longrightarrow\C^{n^2-1}
\]
whose complex differential
\[
\mathrm{d}f_A:T_A SL(n,\C)\longrightarrow\C^{n^2-1}
\]
is injective for every \(A\in SL(n,\C)\). Neither global injectivity nor surjectivity is required.

The source asks whether such a holomorphic immersion exists. For a specified \(n\), the existence condition is
\[
\exists f\in\mathcal O\bigl(SL(n,\C),\C^{n^2-1}\bigr)
\quad\forall A\in SL(n,\C),\qquad
\operatorname{rank}_{\C}(\mathrm{d}f_A)=n^2-1.
\]
The source does not explicitly specify the range of \(n\) or state an outer quantifier over \(n\); these aspects of the intended family of questions require clarification.

The source adds that the requested immersion cannot be obtained algebraically and that the holomorphic existence statement is true for \(n=2\). These are supplied source assertions, not independently verified conclusions. In the usual algebraic interpretation, the first assertion concerns regular algebraic maps with injective differential everywhere, but the source gives neither a definition nor a range of \(n\) for that assertion. In particular, its treatment of the degenerate case \(n=1\) requires clarification. The registry labels the problem open; this supplied label does not override the source's stated positive case \(n=2\) or establish a current literature assessment.
\subsection{Short English statement}
The problem asks whether \(SL(n,\C)\) admits a holomorphic map into complex affine space of the same dimension whose differential is injective everywhere. The source reports a positive answer for \(n=2\) and an obstruction to doing this algebraically. The intended range of \(n\), the scope of the algebraic assertion, and the current status require review.
\subsection{Sources}
\begin{itemize}
\item Local UnsolvedMath registry, permanent ID \(800261\), source code \texttt{catalogue entry}, title \emph{Several Complex Variables Problem List --- Problem 60}. The supplied record labels the problem \texttt{open} and records the year 1995.
\item Buzzard, Fornaess, and Krantz, \emph{The Several Complex Variables Problem List} (1995), Problem 60. Supplied source URL: \url{https://arxiv.org/abs/math/9509202}. The record reports access on 2026-10-16 and extraction by \texttt{arxiv\_latex}; these are supplied provenance details, not a claim of independent source consultation.
\item The problem text attributes the question to R. Narasimhan, Department of Mathematics, University of Chicago, Chicago, IL 60637. The registry's \texttt{proposed\_by} field instead lists Buzzard, Fornaess, and Krantz; both attribution fields are retained, with the latter associated with the source list.
\item The parenthetical assertions concerning algebraic impossibility and the case \(n=2\) are supplied by the problem text. No separate supporting references for those assertions are supplied.
\end{itemize}
\end{document}
